\subsection{Case II: \texorpdfstring{$k > 0$}{k > 0} splitting (1-extension)}
\label{sec:molecular-algebraic-induction-caseII}

\begin{lemma}[Case II rank-lift accounting; the $+D$ core of KT Lemma 6.8]
  \label{lem:case-II-rank-lift}
  \lean{CombinatorialRigidity.Molecular.BodyHingeFramework.rankHypothesis_iff_finrank_pinnedMotions,
        CombinatorialRigidity.Molecular.PanelHingeFramework.rankHypothesis_withNormal_iff_finrank_pinnedMotions,
        CombinatorialRigidity.Molecular.BodyHingeFramework.pinnedMotions_le_withGraph,
        CombinatorialRigidity.Molecular.PanelHingeFramework.toBodyHinge_pinnedMotions_le_withGraph,
        CombinatorialRigidity.Molecular.BodyHingeFramework.pinnedMotions_withGraph_eq,
        CombinatorialRigidity.Molecular.PanelHingeFramework.toBodyHinge_pinnedMotions_withGraph_eq,
        CombinatorialRigidity.Molecular.BodyHingeFramework.hnew_of_isLink_incident,
        CombinatorialRigidity.Molecular.PanelHingeFramework.toBodyHinge_hnew_of_isLink_incident}
  \leanok
  \uses{def:rank-hypothesis, lem:rank-delete-vertex, def:framework-with-graph,
        lem:exists-independent-panel-extensor}
  In the codimension form of \cref{sec:molecular-rigidity-matrix},
  re-inserting (equivalently, pinning) a body $v$ shifts the realization
  count by exactly $D$. A framework $F$ realizes the target
  rank at $k'$ ($\dim Z(G,p) = D + k'$, \cref{def:rank-hypothesis})
  precisely when its body-$v$-pinned motion subspace has dimension $k'$.
  This is the algebraic $+D$ core of the panel-hinge $1$-extension: the
  pinned subspace is the null space of the rigidity matrix with $v$'s $D$
  columns deleted (the smaller graph $G - v$), and $\dim(\text{pinned at
  } v) + D = \dim Z(G,p)$ by the pin-a-body Lemma~5.1
  (\cref{lem:rank-delete-vertex}). Assembled on the panel layer: building
  the $1$-extension by choosing a panel normal $n$ for the re-inserted body
  $v$ (still unhinged, so the choice does not disturb the inductive null
  space, \cref{def:framework-with-graph}), the extended framework
  $P[v \mapsto n]^\flat$ realizes the target rank at $k'$ exactly when the
  splitting-off carries body-$v$-pinned motion dimension $k'$. The graph
  step of the $1$-extension is the unconditional inclusion: re-adding $v$'s
  two new hinge edges (passing from the splitting-off $G_v^{ab}$ up to $G$)
  can only \emph{shrink} the body-$v$-pinned motions, so the inductive
  realization of $G_v^{ab}$ bounds the extended framework's body-$v$-pinned
  dimension from above (the residual cut by the two new edges is the
  genericity-gated Claim~6.9 step). That inclusion is \emph{tight} exactly
  when every base-$v$-pinned motion of $G_v^{ab}$ already clears the two
  re-added $v$-incident hinge constraints $S a \in \operatorname{span}
  C(e_a)$, $S b \in \operatorname{span} C(e_b)$ (a base-pinned motion has
  $S v = 0$): under that hypothesis the body-$v$-pinned motions on $G$ and
  on $G_v^{ab}$ \emph{coincide}, so the extended framework's pinned
  dimension equals the inductive realization's and the $1$-extension lifts
  the rank by exactly $D$. The hypothesis is the honest content of
  Claim~6.9's general position --- supplied by placing the two new
  supporting extensors in general position
  (\cref{lem:exists-independent-panel-extensor}). Since the only links of
  $G$ outside $G_v^{ab}$ are the two $v$-incident edges $v a$, $v b$, and a
  base-pinned motion has $S v = 0$, each re-added hinge constraint
  $S v - S w \in \operatorname{span} C(e)$ collapses to the single
  span-membership $S w \in \operatorname{span} C(e)$ at the non-$v$
  endpoint, reducing the tightness hypothesis to exactly the two conditions
  $S a \in \operatorname{span} C(e_a)$, $S b \in \operatorname{span}
  C(e_b)$ that general position achieves.
\end{lemma}
\begin{proof}
  \leanok
  \uses{lem:rank-delete-vertex, def:framework-with-graph}
  Immediate from the pin-a-body identity
  $\dim(\text{pinned at } v) + D = \dim Z(G,p)$
  (\cref{lem:rank-delete-vertex}): substitute it into the definition of
  the rank hypothesis $\dim Z(G,p) = D + k'$ and cancel the $D$ term. For
  the panel assembly, the body-$v$-pinned motions of $P[v \mapsto n]^\flat$
  agree with those of $P^\flat$ when $v$ is unhinged
  (\cref{def:framework-with-graph}), so the $+D$ accounting applies to the
  pre-choice framework's pinned dimension.
\end{proof}

\begin{lemma}[Edge-substitution bridge for splitting-off]
  \label{lem:splitoff-edge-substitution}
  \lean{Graph.removeVertex_le, Graph.removeVertex_le_splitOff,
        Graph.isLink_incident_of_not_removeVertex}
  \leanok
  The splitting-off $G_v^{ab}$ is \emph{not} a subgraph of $G$: it deletes
  the two $v$-incident edges $v a$, $v b$ but adds a fresh short-circuit
  edge $e_0 \notin E(G)$ joining $a$ and $b$. The two graphs are instead an
  \emph{edge substitution} of one another, sharing the common subgraph
  $G - v$ (all of $G$ away from $v$): $G - v \le G$ and
  $G - v \le G_v^{ab}$. The second inclusion holds because every edge of
  $G - v$ is an edge of $G$, hence $\neq e_0$, so it survives into
  $G_v^{ab}$. Conversely, every link of $G$ \emph{lost} by passing to
  $G - v$ is incident to $v$: the re-added edges over the common lower
  bound are exactly the $v$-incident ones, which is what lets the genericity
  step (\cref{lem:case-II-rank-lift}) collapse each re-added hinge constraint
  on a $v$-pinned motion to a single span membership at its non-$v$ endpoint.
\end{lemma}
\begin{proof}
  \leanok
  Both inclusions are the definition of $\le$ on multigraphs (vertex-set
  and link inclusion). For $G - v \le G_v^{ab}$, a link of $G - v$ is a
  link $e\,x\,y$ of $G$ with $x, y \neq v$; since $e \in E(G)$ and
  $e_0 \notin E(G)$ we have $e \neq e_0$, so $e$ lands in $G_v^{ab}$'s
  $v$-avoiding branch.
\end{proof}

\begin{lemma}[Case II: splitting off a reducible vertex; KT Lemmas 6.7, 6.8]
  \label{lem:case-II}
  \lean{CombinatorialRigidity.Molecular.BodyHingeFramework.isInfinitesimallyRigidOn_insert_iff,
        CombinatorialRigidity.Molecular.PanelHingeFramework.rankHypothesis_withNormal_withGraph_iff_finrank_pinnedMotions}
  \leanok
  \uses{thm:minimal-kdof-reduction, lem:reduction-step,
        lem:rank-delete-vertex, lem:case-II-rank-lift, lem:cycle-realization,
        lem:splitoff-edge-substitution, lem:exists-independent-panel-extensor,
        lem:genericity-device}
  Let $G$ be a minimal $k$-dof-graph with $k > 0$ and a reducible
  degree-2 vertex $v$ (KT~4.6), and let $G_v^{ab}$ be the splitting-off
  at $v$ (a smaller minimal $k$-dof-graph by \cref{lem:reduction-step}).
  A realization of $G_v^{ab}$ at its inductive rank extends to a
  realization of $G$ at rank $D(|V(G)|-1) - k$. This is the panel-hinge
  analogue of Whiteley's bar-joint $1$-extension: re-inserting $v$ with
  its two hinges lifts the rank by exactly $D$ (Lemma~6.8, via the
  Claim~6.9 genericity argument).

  Concretely, with $P$ a panel framework on $G_v^{ab}$ in which $v$ is
  yet unhinged and $P.\mathrm{graph} \le G$, choosing a panel normal $n$
  for $v$ and enlarging the graph to $G$ gives the extended framework
  $(P[v \mapsto n]).\mathrm{withGraph}\,G$, which realizes the target
  rank at $k'$ \emph{iff} the splitting-off $P$ carries body-$v$-pinned
  dimension $k'$ --- the inductive realization lifts to $G$ with the two
  new hinge-row blocks accounting for the $+D$
  (\cref{lem:case-II-rank-lift}). The graph-side hypotheses are
  genericity-free: $v$ is unhinged in $G_v^{ab}$ (true before its two new
  edges are added) and every link of $G$ lost on passing to $G_v^{ab}$ is
  $v$-incident (\cref{lem:splitoff-edge-substitution}). The \emph{one}
  input from the genericity device (Phase~21b, Claim~6.9) is that each
  base-$v$-pinned motion lands in the two new edges' panel-support spans,
  $S a \in \operatorname{span} C(e_a)$, $S b \in \operatorname{span}
  C(e_b)$; this is false pointwise and holds only for the general-position
  normals the rank/dimension count selects
  (\cref{lem:exists-independent-panel-extensor}).
  As with Case~I, the $+D$ lift is accounted $\alpha$-independently through the
  bridge
  $\mathrm{isInfinitesimallyRigidOn\_insert\_iff}$: re-inserting the single
  body $v$ (so $V(G) = \{v\} \cup V(G_v^{ab})$), the framework is rigid on
  $V(G)$ \emph{iff} it is rigid on the splitting-off body set $V(G_v^{ab})$
  (the inductive realization) \emph{and} every motion pins $v$'s screw to the
  bodies it joins, $S v = S w$ for $w \in V(G_v^{ab})$. This references only
  $V(G)$, hence is $\alpha$-independent and composes through the induction. The
  $S v = S w$ condition is the honest geometric content of Claim~6.9's general
  position (the span-membership the device \cref{lem:genericity-device}
  supplies, via \cref{lem:exists-independent-panel-extensor}); the
  $1$-extension count itself ($v$'s two hinges add $D-1$ rows, $v$'s body adds
  one column block) is rank-side and genericity-free. The nullity-side iff is
  retained as the $\alpha$-dependent sibling.
\end{lemma}
\begin{proof}
  \leanok
  \uses{lem:rank-delete-vertex, lem:case-II-rank-lift,
        lem:splitoff-edge-substitution}
  Place $G_v^{ab}$ at its inductive realization; re-insert $v$ joined to
  $a$ and $b$ by hinges chosen generically. The rank-lift accounting
  \cref{lem:case-II-rank-lift} (via the pin-a-body Lemma~5.1,
  \cref{lem:rank-delete-vertex}) accounts for $v$'s screw freedoms through
  the unhinged-$v$ panel-normal invariance, and the genericity-gated
  tightness \cref{lem:case-II-rank-lift} (the $\ge$ half) makes the
  $1$-extension inclusion an equality once the span-membership input
  clears the two new $v$-incident edges (a base-pinned motion has
  $S v = 0$, so each constraint $S v - S w \in \operatorname{span} C(e)$
  collapses to $S w \in \operatorname{span} C(e)$ at the non-$v$
  endpoint). Hence the extended pinned dimension equals the inductive
  one and the rank lifts by exactly $D$, giving $D(|V|-1) - k$. The
  span-membership input is supplied by the genericity device (Claim~6.9,
  Phase~21b).
\end{proof}

\begin{lemma}[Case II realization, all $k \ge 0$: the splitting-off producer; KT Lemma 6.8]
  \label{lem:case-II-realization}
  \lean{CombinatorialRigidity.Molecular.PanelHingeFramework.case_II_realization_all_k}
  \leanok
  \uses{def:panel-hinge-framework, lem:case-II, lem:reduction-step-pos,
        lem:rank-polynomial-of-le-finrank}
  The generic-realization motive is closed under splitting off a
  degree-2 vertex at deficiency $k > 0$. By \cref{lem:reduction-step-pos}
  the splitting-off graph $G_v^{ab}$ is minimal $(k-1)$-dof; the all-$k$
  inductive hypothesis supplies a panel realization $Q$ with
  $\operatorname{rank} R(G_v^{ab}, Q) = D(|V|-2) - (k-1)$.  The
  deficiency-aware eq.-(6.12) placement brick places $v$'s normal at the
  shear $n_a + t\,n_b$ and produces a combined normal assignment $q_0$
  with $N = D(|V|-1) - k$ rows linearly independent in $R(G, q_0)$.  The
  W-suite shear argument (\cref{lem:case-III}) certifies one of the
  candidate $t$-values; the genericity device
  (\cref{lem:rank-polynomial-of-le-finrank}) re-realizes the attained rank
  at a general-position, algebraically-independent seed.
\end{lemma}
\begin{proof}
  \leanok
  \uses{lem:case-III, lem:rank-polynomial-of-le-finrank}
  The deficiency-aware placement brick (analogous to
  \cref{lem:case-II-realization-placement} at $k > 0$) supplies $N \le
  \operatorname{rank} R(G, q_0)$.  The W-suite (\cref{lem:case-III}) finds
  a good shear parameter; \cref{lem:rank-polynomial-of-le-finrank} lifts
  the rank lower bound to a nonzero rational polynomial whose non-roots
  realize deficiency $k$.
\end{proof}

